\section{Beschreibung des Unternehmens}

Brown \& Brown, Inc. ist ein b\"orsennotierter Versicherungsmakler mit Sitz in
Daytona Beach, Florida. Das Unternehmen vermittelt Versicherungsschutz und
erbringt Risikodienstleistungen; es tr\"agt das Versicherungsrisiko nur in
eng begrenzten Ausnahmen selbst und verdient daher \"uberwiegend Provisionen
und Geb\"uhren. Zum 31.\,Dezember 2025 besch\"aftigte der Konzern
\num{\Mitarbeiter} Personen weltweit.\quelleGB{2025}{7}\quelleMerken{s3gb}
Seit dem dritten Quartal 2025 berichtet das Unternehmen in zwei Segmenten:
Retail mit \num{\MitarbeiterRetail} Besch\"aftigten\quelleGB{2025}{6}\quelleMerken{s3gb} und
Specialty Distribution mit \num{\MitarbeiterSpez}.\quelleGB{2025}{10}\quelleMerken{s3gb}
Die Aktie ist an der New York Stock Exchange unter dem K\"urzel BRO
notiert.\quelleGB{2025}{28}\quelleMerken{s3gb}

\subsection{Handelnde Personen in f\"uhrender Position}

Die Gesch\"aftsf\"uhrung liegt bei \num{\VorstandsmitgliederAnzahl}
executive officers.\quelleErneut{s3gb} An ihrer Spitze
steht J.\,Powell Brown als president und chief executive officer; er f\"uhrt
das Unternehmen seit Juli 2009 und geh\"ort dem Board seit Oktober 2007 an.
Vorsitzender des Board of Directors ist sein Vater J.\,Hyatt Brown, der das
Unternehmen von 1993 bis 2009 selbst geleitet hat. Ein dritter Angeh\"origer
derselben Familie, P.\,Barrett Brown, ist executive vice
president.\quelleErneut{s3gb} Finanzvorstand und treasurer ist
R.\,Andrew Watts; die beiden Segmente werden von Stephen P.\,Hearn (Retail,
zugleich chief operating officer) und Stephen M.\,Boyd (Specialty
Distribution) gef\"uhrt. Die Zuk\"aufe verantwortet mit J.\,Scott Penny ein
eigener chief acquisitions officer.\quelleErneut{s3gb}

Das Board of Directors umfasst \num{\AufsichtsratMitglieder}
Mitglieder.\quelleQM{Proxy Statement 2026}{76}\quelleMerken{s3qmproxystateme}
\bk{Bemerkenswert ist dabei weniger die Zahl als die Verschr\"ankung: Der
Vorsitzende des Kontrollgremiums ist zugleich der gr\"o{\ss}te
Einzelaktion\"ar und der Vater des Vorstandsvorsitzenden. Die Aufteilung der
Rollen zwischen Vorsitz und Gesch\"aftsf\"uhrung besteht damit formal, nicht
wirtschaftlich. Dass eigens ein chief acquisitions officer im Rang eines
executive vice president gef\"uhrt wird, ist die organisatorische Antwort
auf ein Gesch\"aftsmodell, dessen Wachstum aus Zuk\"aufen stammt.}

\subsection{Eigent\"umerstruktur}

Tabelle~\ref{tab:eigentuemer} gibt die zum Stichtag 2.\,M\"arz 2026 im Proxy
Statement ausgewiesenen Anteile wieder.

\begin{table}[htbp]
\centering
\caption{Anteilseigner mit ausgewiesener Beteiligung, Stichtag 2.~M\"arz 2026.
Die Prozentwerte sind die des Proxy Statement.}
\label{tab:eigentuemer}
\begin{tabular}{lrr}
\toprule
Anteilseigner & St\"uck & Anteil in \%\\
\midrule
J.~Hyatt Brown (Vorsitzender)        & \num{\HyattBrownStueck}   & \num{\HyattBrownAnteil}\\
The Vanguard Group, Inc.             & \num{\VanguardStueck}     & \num{\VanguardAnteil}\\
Capital World Investors              & \num{\CapitalWorldStueck} & \num{\CapitalWorldAnteil}\\
J.~Powell Brown (Vorstandsvorsitzender) & \num{\PowellBrownStueck} & \num{\PowellBrownAnteil}\\
\midrule
Organe insgesamt (\num{\OrganePersonen} Personen) & \num{\OrganeStueck} & \num{\OrganeAnteil}\\
\bottomrule
\end{tabular}
\end{table}

\quelleQM{Proxy Statement 2026}{76}\quelleMerken{s4qmproxystateme}Die beiden zuletzt genannten Personen sind in der
Zeile der Organe bereits enthalten; die Spalte summiert sich deshalb
absichtlich nicht. Ein beherrschender Aktion\"ar besteht nicht: Der gr\"o{\ss}te
Einzelbestand liegt bei \Pct{\HyattBrownAnteil}, die Organe zusammen halten
\Pct{\OrganeAnteil}. Der \"ubrige Bestand ist Streubesitz, in dem mit
Vanguard und Capital World zwei Verm\"ogensverwalter mit ausgewiesenen
Anteilen von \Pct{\VanguardAnteil} beziehungsweise \Pct{\CapitalWorldAnteil}
stehen. Weitere Halter oberhalb von f\"unf Prozent nennt das Proxy Statement
nicht.\quelleErneut{s4qmproxystateme}

Die Zahl der umlaufenden Aktien ist 2025 sprunghaft gestiegen, von
\num{\AktienUmlaufVJ} auf \num{\AktienUmlauf} Millionen
St\"uck.\quelleGB{2025}{51}\quelleMerken{s4gb} Ursache ist die
Finanzierung der Accession-\"Ubernahme (Abschnitt~\ref{sec:accession}). Zum
10.\,Februar 2026 waren \num{\AktienRegisterStueck} Aktien im Umlauf,
gehalten von rund \num{\AktionaereRegister} eingetragenen
Aktion\"aren.\quelleGB{2025}{28}\quelleMerken{s4gb}

\subsection{Beurteilung durch die Kapitalm\"arkte}

Zum 31.\,Dezember 2025 schloss die Aktie bei \USDje{\KursSchlussZFV}; bei
\num{\AktienUmlauf} Millionen umlaufenden Aktien entspricht das einer
Marktkapitalisierung von \MioUSD{\Marktkapitalisierung}. Beide Gr\"o{\ss}en
tragen dasselbe Stichtagsdatum. Ein Jahr zuvor lag der Wert bei
\MioUSD{\MarktkapitalisierungVJ}: Der Markt bewertete das Unternehmen Ende
2025 also niedriger als Ende 2024, obwohl der Umsatz im selben Zeitraum um
\Pct{\UmsatzDelta} gewachsen ist. Auf das bilanzielle Eigenkapital bezogen
betr\"agt die Bewertung das \Faktor{\KursBuchwert}-fache, auf das
verw\"asserte Ergebnis je Aktie das \Faktor{\KGV}-fache nach
\Faktor{\KGVVJ} im Vorjahr.

\subsection{Analystenabdeckung}

Die Investor-Relations-Seite des Unternehmens nennt
\num{\AnalystenHaeuser} abdeckende H\"auser, jeweils mit dem Namen des
zust\"andigen Analysten. Vertreten sind unter anderem Barclays, BofA
Securities, BMO Capital Markets, Citi, Goldman Sachs, J.\,P.\,Morgan,
Jefferies, Keefe, Bruyette \& Woods, Morgan Stanley, Raymond James, UBS und
Wells Fargo.\quelleWeb{Brown \& Brown, Inc., Analyst Coverage}{quelle:analysten}{28.08.2026}\quelleMerken{s4webbrownbrowni}

\subsection{Botschaften und L\"ucken}
\label{sec:luecken}

Die Liste der abdeckenden H\"auser ist die einzige Angabe, die das
Unternehmen zur Analystenmeinung ver\"offentlicht. Weder Empfehlungen noch
Kursziele werden genannt; die Seite h\"alt ausdr\"ucklich fest, dass
Einsch\"atzungen und Prognosen der Analysten allein deren eigene sind und
nicht die des Unternehmens
wiedergeben.\quelleErneut{s4webbrownbrowni} Diese Analyse
gibt deshalb weder ein Konsensrating noch ein durchschnittliches Kursziel
an. Die Beurteilung des Kursverlaufs in Teil~\ref{sec:kurs} st\"utzt sich
statt dessen auf eine Gr\"o{\ss}e, die das Unternehmen selbst quartalsweise
ver\"offentlicht: das organische Umsatzwachstum.

Zwei weitere L\"ucken sind zu benennen. Erstens ver\"offentlicht das
Unternehmen als US-amerikanischer Emittent keine Jahresprognose in Form
einer Umsatz- oder Margenbandbreite; ein Abgleich von Ausblick und Ist, wie
ihn ein deutscher Gesch\"aftsbericht erlaubt, ist deshalb nicht m\"oglich
(Abschnitt~\ref{sec:ausblick}). Zweitens weist das Proxy Statement neben
Vanguard und Capital World keinen weiteren Halter oberhalb von f\"unf
Prozent aus; ob andere institutionelle Investoren knapp darunter liegen,
l\"asst sich aus der Prim\"arquelle nicht entnehmen und wird hier nicht
gesch\"atzt.
